\documentclass[11pt]{article}

% ---- Packages ---------------------------------------------------------------
\usepackage[margin=1in]{geometry}
\usepackage{amsmath}               % equations
\usepackage{booktabs}              % clean tables: \toprule, \midrule, \bottomrule
\usepackage{graphicx}              % figures
\usepackage[authoryear,round]{natbib}  % \citet{key} -> Corsi (2009); \citep{key} -> (Corsi, 2009)
\usepackage[hidelinks]{hyperref}   % clickable references; load last

\graphicspath{{figures/}}          % upload the PDFs from output/ into a figures/ folder

% ---- Title ------------------------------------------------------------------
\title{Do Better Volatility Forecasts Improve Avellaneda-Stoikov Market Making?\\
       Evidence from Bitcoin Perpetual Futures}
\author{Author Name\\ University}
\date{\today}

\begin{document}
\maketitle

\begin{abstract}
% Write last. About 150 words: question, data, method, main result, why it matters.
\end{abstract}

% ============================================================================
\section{Introduction}
\label{sec:intro}
% Write second to last. The question, why it matters, the answer, the contributions.

% ============================================================================
\section{Related literature}
\label{sec:literature}
% AS and its extensions; volatility forecasting (HAR, GARCH, QLIKE);
% adverse selection. See paper/literature_notes.md for what each source supports.

% ============================================================================
\section{Data}
\label{sec:data}
The data come from Binance's public market-data archive\footnote{\url{https://data.binance.vision}}
for the BTCUSDT perpetual futures contract. Two feeds are used. The \texttt{bookTicker} feed records
every change to the best bid and ask, with millisecond timestamps. From it we build the mid-price,
$s_t = (b_t + a_t)/2$, where $b_t$ and $a_t$ are the best bid and ask at time $t$. The mid-price drives both the volatility estimators and the market maker's quotes. The
\texttt{aggTrades} feed records every trade with its price, size and aggressor side. The simulator uses
these trades to decide when the market maker's quotes are filled (Section~\ref{sec:simulator}).

The archive has two limitations that shape the study. First, it contains only the best bid and ask
(level~1), with no depth behind them. The simulator therefore cannot know the market maker's place in the
queue at a price level, so fills are bounded by a conservative and an optimistic rule rather than modeled
exactly (Section~\ref{sec:simulator}). Second, the futures \texttt{bookTicker} feed is published only from
16~May~2023 to 30~March~2024. Trades are available over a much longer period, but a market maker cannot
be simulated without quotes, so this window bounds the study.

\paragraph{Sample.}
Table~\ref{tab:sample} shows how the sample is built. It starts from every weekday from 17~May~2023 to
29~March~2024 (228 days); the feed's first day is dropped in case it is partial. Weekends are excluded
because their trading volume is much lower: including them would mix a change in liquidity into a
comparison meant to isolate volatility.

Seven days are excluded for broken quote data. The archive has holes: stretches with no quote updates
while trading continues. During such a gap the simulator's quotes stay at stale prices and are filled by
trades at prices the real market has already left. On the affected days the market maker built positions
of up to about 100~BTC, against a quote size of 0.005~BTC. Every day is therefore screened: a day fails if
it contains a quote gap of 30 seconds or more with trades inside it, or if quotes are missing at the start
or end of the day. Five days fail the screen, with gaps of 8 to 105 minutes. Two more, whose quote feed
starts several hours late, were excluded by hand before the screen existed.

Two development days, 19~July~2023 and 10~January~2024, are also excluded. The pipeline was built and its
settings tuned on them, so keeping them would test the method on the data it was fitted to. This leaves
219 study days.

\paragraph{Hold-out.}
Every design decision was made on an earlier sample of 41 days, which included the two development days.
The 180 study days outside that sample were added only after every decision was final, so they form an
out-of-sample hold-out (Section~\ref{sec:robustness}).

\begin{table}[ht]
  \centering
  \caption{Sample construction.}
  \label{tab:sample}
  \begin{tabular}{lr}
    \toprule
    Step & Days \\
    \midrule
    Weekdays, 17 May 2023 -- 29 March 2024 & 228 \\
    Excluded: broken quote data & $-7$ \\
    Excluded: development days & $-2$ \\
    \midrule
    Study sample & 219 \\
    \quad of which out-of-sample hold-out & 180 \\
    \bottomrule
  \end{tabular}
\end{table}

% ============================================================================
\section{Method}
\label{sec:method}

\subsection{Volatility estimators}
\label{sec:estimators}

We compare four estimators that span the main approaches to short-horizon volatility: a model-free
measure (realized volatility), exponential smoothing (EWMA), a fitted conditional-variance model
(GARCH(1,1)) and a regression on realized variances (HAR-RV). All four are built from one-second log
returns of the mid-price, $r_t = \ln(s_t / s_{t-1})$, where $s_t$ here is the mid-price at the end of
second $t$ (the last quote update in that second, or the previous second's value if there is none). Each
estimator produces a per-second variance $h_t$.
Because the Avellaneda-Stoikov model describes the price in dollars, every estimator's output is
converted to dollar volatility in the same way, $\sigma_t = s_t\sqrt{h_t}$ (in USD per
$\sqrt{\text{second}}$), so the estimators differ only in how they estimate $h_t$. All are causal: the
value in force at time $t$ uses only data observed before $t$, and every refit uses only data before the
refit time.

\paragraph{Realized volatility (RV).}
The mean of squared returns over the last $W = 300$ seconds,
\begin{equation}
  h_t^{RV} = \frac{1}{W}\sum_{i=0}^{W-1} r_{t-i}^2 .
  \label{eq:rv}
\end{equation}
Realized variance is the standard model-free volatility measure \citep{andersen1998}; in a comparison
of over 400 realized measures, few significantly beat plain five-minute realized variance
\citep{lps2015}. It is also how the best-known open-source
Avellaneda-Stoikov implementation estimates $\sigma$ \citep{hummingbot_instantvol}, which makes it the
natural baseline. The 300-second window is our choice: it matches the five-minute horizon over which
the market maker uses $\sigma$ (Section~\ref{sec:as}).

\paragraph{EWMA.}
The RiskMetrics recursion \citep{riskmetrics1996},
\begin{equation}
  h_t^{EWMA} = \lambda h_{t-1}^{EWMA} + (1-\lambda)\, r_t^2 .
  \label{eq:ewma}
\end{equation}
The usual $\lambda = 0.94$ is calibrated for daily returns and has no meaning at one-second sampling.
We instead set $\lambda$ so that the average age of the data EWMA weights, $\lambda/(1-\lambda)$, equals
that of RV's window, $(W-1)/2 = 149.5$ seconds, which gives $\lambda \approx 0.9934$. RV and EWMA then
use equally recent information and differ only in how they weight it.

\paragraph{GARCH(1,1).}
The conditional variance follows \citet{bollerslev1986},
\begin{equation}
  h_t^{G} = \omega + \alpha\, r_{t-1}^2 + \beta\, h_{t-1}^{G} ,
  \label{eq:garch}
\end{equation}
where $\alpha$ measures the response to the latest squared return, $\beta$ the carry-over of the previous
variance, and $\omega$ pulls the variance toward its long-run level $\omega/(1-\alpha-\beta)$. The model
captures volatility clustering, and \citet{hansenlunde2005} find it hard to beat among 330 ARCH-type
models. Parameters are estimated by
maximum likelihood (zero mean, normal errors) on the trailing three hours of returns and re-estimated
every hour; between refits the recursion runs every second. Fewer than 1\% of the hourly fits fail to
converge, in which case the previous fit is kept. The fitted persistence $\alpha + \beta$ has a median
of 0.98, so a volatility shock decays by half in about 34 seconds. Since EWMA is the special case $\omega = 0$, $\alpha + \beta = 1$, comparing the two tests
whether fitting the parameters adds anything.

\paragraph{HAR-RV.}
\citet{corsi2009} forecasts realized variance from its averages over several past horizons. We adapt
his daily, weekly and monthly components to intraday horizons of 1, 5 and 30 minutes:
\begin{equation}
  h_t^{HAR} = \beta_0 + \beta_1 RV_t^{(1\,\text{min})} + \beta_2 RV_t^{(5\,\text{min})}
              + \beta_3 RV_t^{(30\,\text{min})} ,
  \label{eq:har}
\end{equation}
where $RV_t^{(w)}$ is the mean squared return over the trailing window $w$ and $\beta_0,\dots,\beta_3$
are regression coefficients. Unlike the other three
estimators, HAR is fitted directly to the quantity the market maker needs: the mean squared return over
the next 300 seconds. The coefficients are estimated by non-negative least squares on all of the day
so far, sampled every 10 seconds, and re-estimated every hour. An expanding window is used because a
trailing three-hour window holds only about six independent observations of the 30-minute component,
which produced negative coefficients and negative variance forecasts in development. Forecasts are
floored at 1\% of the day's mean squared return, a floor that binds for only a few seconds in the whole
sample. HAR is included because it beats GARCH-type models in Bitcoin volatility forecasting
\citep{aganin2023,bergsli2022}.

\paragraph{Raw volatility and burn-in.}
Each estimator's $\sigma_t$ enters the quoting formula unadjusted, so any bias in its level counts as
part of that estimator's quality. A variant rescaled every hour to the trailing three-hour realized
variance is reported as a robustness check (Section~\ref{sec:robustness}). The first three hours of each
day (00:00--03:00 UTC) serve as burn-in for the first fits; performance is scored from 03:00 UTC.

\paragraph{Microstructure noise.}
At very short sampling intervals, realized volatility is usually biased upward by bid-ask bounce
\citep{bandirussell2008}. On Binance mid-quotes the opposite happens (Figure~\ref{fig:signature}), consistent with evidence that
noise in quote data is not simple bid-ask bounce \citep{hansenlunde2006}: on
a calm day 79\% of one-second returns are exactly zero, and one-second realized volatility sits about
17\% below its peak at 15--20-second sampling, against about 6\% on a volatile day. Because this bias
varies with the volatility regime, quantities used to compare hours, such as the regime labels
(Section~\ref{sec:evaluation}), use 60-second returns instead. The estimators keep one-second returns,
the sampling a live market maker would use, and their raw output carries this bias, as noted above.

\begin{figure}[ht]
  \centering
  \includegraphics[width=\textwidth]{signature_plot.pdf}
  \caption{Volatility signature plot for the BTCUSDT mid-price on the two development days: annualized
  realized volatility at each sampling interval, divided by its value at 300 seconds. Realized volatility
  is not inflated at short intervals; on the calm day (19 July 2023) it is biased downward instead.}
  \label{fig:signature}
\end{figure}

\subsection{The Avellaneda-Stoikov model}
\label{sec:as}

The market maker quotes with the model of \citet{avellaneda2008}. A risk-averse market maker posts a
bid and an ask around a \emph{reservation price}, its own valuation of the asset given the inventory it
holds,
\begin{equation}
  r_t = s_t - q_t\, \gamma\, \sigma_t^2\, \tau_t ,
  \label{eq:reservation}
\end{equation}
and sets the total spread between its two quotes to
\begin{equation}
  \delta^a_t + \delta^b_t = \gamma\, \sigma_t^2\, \tau_t + \frac{2}{\gamma}\ln\left(1 + \frac{\gamma}{k_t}\right),
  \label{eq:spread}
\end{equation}
placed symmetrically around $r_t$. Here $s_t$ is the mid-price, $q_t$ the signed inventory in BTC,
$\gamma$ the risk aversion, $\sigma_t$ the volatility from the estimator under test, $\tau_t$ the time
remaining in the horizon, and $k_t$ the order-flow parameter defined below. The two equations play
different roles. Inventory enters only the reservation price: a long position lowers $r_t$ below the
mid, shifting both quotes down so that the market maker sells more readily and sheds risk. The width of
the spread does not depend on inventory; it has a volatility term, $\gamma\sigma_t^2\tau_t$, and an
order-flow term set by $k_t$. These closed-form quotes are the approximate solution that
\citeauthor{avellaneda2008} derive for their model, and they are the form used in practice.

\paragraph{Horizon.}
The model assumes a finite trading horizon, with $\tau_t$ counting down to its end. We use a horizon of
300 seconds: $\tau_t$ counts down from 300 to 0 and then restarts. A full trading day would make the
volatility term unusable. For an illustrative $\sigma_t = \$4$ per $\sqrt{\text{second}}$ and
$\gamma = 10^{-3}$, that term is about \$5 over 300 seconds but about \$1{,}400 over a day, far wider
than any quote that would trade. The 300-second horizon also matches the RV window
(Section~\ref{sec:estimators}).

\paragraph{Order flow.}
\citet{avellaneda2008} model fills as arriving at rate $\lambda(\delta) = A e^{-k\delta}$, so a quote
$\delta$ dollars from the mid is filled less often the farther away it is; only $k$ enters the quotes.
Following their own calibration idea, we estimate $k$ from the data. For each 100-millisecond decision
window we record how far the market reached from the mid on each side, through a trade or the opposite
best quote. This gives the rate $\lambda(\delta)$ at which a quote at each distance from \$0.10 to \$12
would have been reached, and $k$ is minus the slope of a least-squares fit of $\ln\lambda(\delta)$ on
$\delta$. The fit uses the trailing three hours and is repeated every hour, balancing enough data for a
stable estimate against adapting to changing conditions. The fits are close ($R^2$ of 0.98 to 0.995), and
$k$ ranges from about 0.2 per dollar in volatile hours to about 1.2 in calm ones. Fitting $k$ once per
day or holding it constant are robustness checks (Section~\ref{sec:robustness}).

\paragraph{Risk aversion and order size.}
The risk aversion $\gamma$ is a preference, not a quantity that can be estimated from data, and choosing
the best value after seeing the results would be selection on outcomes. Results are therefore reported
for four values, $\gamma \in \{10^{-5}, 10^{-4}, 10^{-3}, 10^{-2}\}$. Across this range the volatility
term's share of the spread rises from a few percent to most of it (Section~\ref{sec:results}), which
governs how much the choice of estimator can matter. Each order is for a fixed 0.005~BTC (about \$120 to
\$370 over the sample). The simulator assumes the market maker's orders do not move the market, which is
plausible only for small orders. Inventory is not capped; the reservation price alone controls it.

\subsection{Simulator}
\label{sec:simulator}

The market maker is simulated on the raw quote and trade events, replayed in time order. Each day is
simulated independently, starting with zero inventory, and the simulation has no random elements: the
same inputs always produce the same fills.

\paragraph{Decisions and orders.}
Every 100 milliseconds the market maker recomputes its quotes from Equations~\eqref{eq:reservation}
and~\eqref{eq:spread}, using only information available before the decision: the latest mid-price,
the volatility estimate from the last completed second, the current $k_t$ and its inventory. Quotes are
rounded to the \$0.10 tick, the bid down and the ask up, and reach the exchange 50 milliseconds later.
Orders are post-only: a quote that would cross the book when it arrives is rejected rather than executed
as a taker order. Each quote rests from its arrival until the next quote arrives, and only events in
that window can fill it.

\paragraph{Fills.}
With only the best bid and ask available, the market maker's place in the queue at its price level is
unknown, so fills are bounded by two deterministic rules. Under the \emph{conservative} rule, a quote is
filled only when the market trades through it: a trade prints at a strictly better price, or the
opposite best quote moves to or past it. Either event means the whole price level, including our order,
was taken. The \emph{optimistic} rule also fills a quote from trades exactly at its price, up to their
size, as if the market maker were first in the queue. Any queue model would fall between the two. All
fills are at the quote price. A fill through the quote is always for the full 0.005~BTC; a fill from
trades exactly at the quote price, which only the optimistic rule allows, is for the size of those trades,
up to 0.005~BTC.
In practice the two rules give nearly the same results, within 1--2\% in P\&L, because the quotes sit
a few dollars from the mid, tens of ticks away in a market whose spread is usually a single tick. At that
distance fills come from trades sweeping through the level rather than from queueing at it.

\paragraph{P\&L and fees.}
P\&L is cash plus inventory marked to the mid-price at the end of every second. The objective in
\citet{avellaneda2008} contains no transaction costs, so the quotes do not depend on fees, and neither do
the fills. Net P\&L is therefore exactly gross P\&L minus the fee rate times the traded notional, and fees
are applied after the simulation, at 0 and 2 basis points per side (roughly Binance's standard maker
fee).

\paragraph{Validation.}
The fill rules and the P\&L decomposition (Section~\ref{sec:benchmarks}) are tested on hand-built
scenarios with known outcomes: a trade exactly at the quote, a trade through it, the opposite quote
crossing it and a post-only rejection, each under both rules.

\subsection{Benchmarks and P\&L decomposition}
\label{sec:benchmarks}

Two additions help explain the comparison: benchmark volatility series that bracket the estimators,
and a decomposition of each run's P\&L into its sources.

\paragraph{Benchmarks.}
Three benchmark series run through the same simulator but are not estimators: two of them use future
data and could not be traded, and all three are kept out of the estimator comparisons and their tests.
Each is converted to dollars in the same way as the estimators, $\sigma_t = s_t\sqrt{h_t}$.
\begin{itemize}
  \item \emph{Flat}: the realized variance of the trailing three hours, updated every hour and held
    constant within it. It is tradable and carries the volatility level with no dynamics within the hour.
  \item \emph{OracleHour}: the realized variance of the current clock hour, from that hour's own
    returns, held constant within it. It knows the hour's true level but not when within the hour
    volatility arrives.
  \item \emph{Oracle}: the realized variance of the next 300 seconds, the same target the forecast
    loss scores against (Section~\ref{sec:evaluation}). It is the perfect volatility forecast for the
    market maker's horizon.
\end{itemize}
The gaps between them separate what volatility information is worth. Flat to OracleHour is the value of
knowing the hour's level; OracleHour to Oracle is the value of also knowing its timing within the hour.
Oracle minus RV bounds how much any better volatility forecast could improve on RV within this quoting
rule.

\paragraph{P\&L decomposition.}
A market maker earns the spread on each fill but loses when prices move against the inventory it
takes on, and the loss right after a fill reflects trading with better-informed counterparties, known as
adverse selection \citep{glostenmilgrom1985}. We decompose P\&L following \citet{hendershott2011}, who
split the effective spread on each trade into a realized spread, the part the liquidity supplier keeps,
and a price impact, the part lost to the subsequent move of the quote midpoint. For a fill of $x$ BTC in
direction $d$ ($+1$ for a buy, $-1$ for a sell) at price $p$, quoted when the mid was $m_0$,
\begin{align}
  \text{spread captured} &= x\, d\, (m_0 - p), \label{eq:spread_captured}\\
  \text{adverse selection} &= x\, d\, (m_H - m_0), \label{eq:adverse_selection}
\end{align}
where $m_H$ is the mid at the end of the $H$-th second after the fill. Spread captured is the fill's
edge over the mid its quote was based on, and is positive unless the inventory skew has pushed the quote
past the mid. Adverse selection is the mid's move in the fill's direction over the next $H$ seconds, and
is negative when the price moves against the market maker. Their sum is the realized spread. Spread
captured corresponds to the effective spread of \citet{hendershott2011} and adverse selection to minus
their price impact, both seen from the market maker's side.

Three adaptations are ours. The reference is the mid the quote was based on rather than the mid at the
time of the fill, so price moves during the order's latency count as adverse selection. The horizon is
$H = 5$ seconds, with 60 seconds as a check, rather than minutes, matching how quickly quotes change in
this market. And because the decomposition of \citet{hendershott2011} applies to each trade rather than
to a trader's P\&L, we add a third component: everything else, the price moves on positions held longer
than $H$ seconds, is the \emph{inventory} component. The three components add up exactly to the P\&L in
every second; fees are a fourth, separate component. Comparing the components across estimators shows
whether an estimator changes P\&L through the spread it earns, through how often it is picked off, or
through the inventory it carries.

\subsection{Evaluation}
\label{sec:evaluation}

% ============================================================================
\section{Results}
\label{sec:results}
% Economics, forecast accuracy, P&L comparison, power and equivalence,
% decomposition, sigma's role in the quotes, benchmarks.
%
% Figure example (delete or replace):
% \begin{figure}[ht]
%   \centering
%   \includegraphics[width=\textwidth]{regime_comparison_raw_conservative_fee0.pdf}
%   \caption{...}
%   \label{fig:regime}
% \end{figure}
%
% Table example (delete or replace):
% \begin{table}[ht]
%   \centering
%   \caption{...}
%   \label{tab:example}
%   \begin{tabular}{lrr}
%     \toprule
%     Estimator & Column 1 & Column 2 \\
%     \midrule
%     RV & 0.00 & 0.00 \\
%     \bottomrule
%   \end{tabular}
% \end{table}

% ============================================================================
\section{Robustness and out-of-sample evidence}
\label{sec:robustness}

% ============================================================================
\section{Conclusion}
\label{sec:conclusion}
% Findings, limitations, future work.

% ---- References -------------------------------------------------------------
\bibliographystyle{plainnat}
\bibliography{references}

\end{document}
